\begin{result}{F.3.29}{A single normalizer line in genus three}{Proved; every prime and unrestricted Jacobian}
Let $V$ be a six-dimensional symplectic vector space over $\Q_\ell$.
Let $U_{3,6}\subset\mathfrak h_{3,6}$ and
$q_6:U_{3,6}\to Q_6$ be the actual universal arithmetic image and
the one-dimensional quotient in F.3.24 and Nakamura's original
(4.7.1)--(4.7.2), printed p.~117.
For every nonzero $a\in V$, put
\[
 P_a^{\rm all}=\{D\in\mathfrak h_{3,6}:
                         D(a)\in[\mathbb L_6(V),a]\}.
\]
Then
\[
                    q_6(U_{3,6}\cap P_a^{\rm all})=0.       \tag{SL.1}
\]
Only this one line is constrained. In particular, no conjugate line,
CM field, multiplicity space, or nondegenerate span of Galois conjugates
is an assumption of (SL.1).

For every smooth proper genus-three curve $C/\Q$ with a rational
point $y$, every prime $\ell$, and every finite extension $K/\Q$,
a nontrivial closed procyclic subgroup of
$\pi_1((C\setminus\{y\})_{\overline K})^{(\ell)}$
whose arithmetic normalizer has open image is contained in actual
cusp inertia. The Jacobian is unrestricted.
\end{result}
\begin{proof}
We use the actual mixed geometric-normalizer equations on the original
six-dimensional space and an exact contraction certificate.

\paragraph{The actual image maps.}
Choose a symplectic basis $a,b,c,d,e,f$, with pairs $(a,b),(c,d),(e,f)$.
Write $J_m$ for the actual geometric Johnson image in special
degree-$m$ derivations. The strict geometric-action construction of
F.3.24 gives
\[
 J_6\subset U_{3,6},\qquad q_6(J_6)=0,\qquad
                  [D,J_1]\subset J_7\quad(D\in U_{3,6}).   \tag{SL.2}
\]
These maps use Hain's strict weight filtration and the actual arithmetic
normalization of the geometric monodromy image, as checked in F.3.24.
The first inclusion and quotient are also explicit in Nakamura's
(4.7.1) and the last paragraph of printed p.~117
\cite{AuditNakamura1996}. Thus the brackets in (SL.2) are vectors in
the geometric image. The argument does not identify a formal
closed-surface ideal with that image.

For $u,v,w\in V$, set
\[
 \delta_{uvw}(x)=\langle x,u\rangle[v,w]
              +\langle x,v\rangle[w,u]
              +\langle x,w\rangle[u,v].
\]
These derivations give $J_1=\Lambda^3V$. In particular, the twenty
triples of distinct basis letters span $J_1$.

\paragraph{One additional primary-source detector.}
The Enomoto--Satoh map vanishes on $J_7$ by the source maps already
verified in F.3.24. The newly opened primary preprint
\cite[Propositions~4.1--4.2, p.~11]{AuditMoritaSakasaiSuzuki2026v1}
gives an additional symplectic-equivariant map
$R:V^{\otimes9}\to V$ whose restriction to $J_7$ is zero.
Explicitly,
\[
\begin{split}
 R(u_1\otimes\cdots\otimes u_9)
  =\bigl(&-2\langle u_1,u_2\rangle\langle u_3,u_5\rangle
                \langle u_4,u_7\rangle\langle u_6,u_8\rangle\\
         &+2\langle u_1,u_2\rangle\langle u_3,u_6\rangle
                \langle u_4,u_7\rangle\langle u_5,u_8\rangle\\
         &+2\langle u_1,u_2\rangle\langle u_3,u_6\rangle
                \langle u_4,u_8\rangle\langle u_5,u_7\rangle\\
         &+2\langle u_1,u_2\rangle\langle u_3,u_7\rangle
                \langle u_4,u_6\rangle\langle u_5,u_8\rangle\\
         &+\langle u_1,u_3\rangle\langle u_2,u_6\rangle
                \langle u_4,u_7\rangle\langle u_5,u_8\rangle\bigr)u_9.
\end{split}                                                   \tag{SL.3}
\]
The tensor map for a degree-seven derivation is
$D\mapsto\sum_i(x_i\otimes D(y_i)-y_i\otimes D(x_i))$.
For the genus-three use, choose an embedding of the boundary surface
in a surface of genus $G$ in the source's stable range. Extending a
mapping class by the identity gives the actual geometric map on
Johnson images, and its derivation extends by zero on the added
handles. For the symplectic inclusion $i:V\hookrightarrow V_G$,
every contraction in (SL.3) gives
$R_G(i^{\otimes9}T)=iR_3(T)$. The source's vanishing on the
extended actual Johnson image and the injectivity of $i$ therefore
give $R_3(J_7)=0$. No stable multiplicity count is imposed on
the original genus-three space.
Consequently every actual $D\in U_{3,6}$ satisfies
\[
 \operatorname{ES}_7([D,\delta])=0\quad(\delta\in J_1),
 \qquad [c]R([D,\delta_{abc}])=0,                         \tag{SL.4}
\]
where $[c]$ denotes the coefficient of $c$ in $V$.
Only this one scalar component of $R$ is used. In particular, the
final proof does not impose a raw exterior-power equation in place
of the quotient target of the separate map $S$ in that preprint.

The source gives a sketch of proof for Proposition~4.1. Our certificate
copies each of the five displayed contractions, and independently
checks its vanishing on the forty actual degree-seven brackets
$[Z_i,\delta_{uvw}]$ used below. The universal vanishing on $J_7$
is the stated primary-source result; the forty checks are a separate
test of the implemented formula and conventions.

\paragraph{The complete balanced one-line space.}
Use the diagonal similitude torus with weights
$a:\alpha_1,b:t/\alpha_1,c:\alpha_2,d:t/\alpha_2,
e:\alpha_3,f:t/\alpha_3$. Let $P_a$ be the $t^3$ part of
$P_a^{\rm all}$. Its specialness output has multidegree
$(n_1,n_1,n_2,n_2,n_3,n_3)$, where $n_1+n_2+n_3=4$.
For $D(a)$ use the entire image $[\mathbb L_6,a]$ in the
required multidegree. For every other generator, including $b$,
use the entire degree-seven Lyndon basis in its required multidegree.
The specialness equations have the following exact ranks:
\[
\begin{array}{c|r|r|r|r}
n&\text{rows}&\text{columns}&\text{rank}&\text{nullity}\\\hline
(0,0,4),(0,4,0)&8&10&8&2\\
(0,1,3),(0,3,1)&140&160&140&20\\
(0,2,2)&312&360&312&48\\
(1,0,3),(1,3,0)&140&143&140&3\\
(1,1,2),(1,2,1)&1260&1290&1260&30\\
(2,0,2),(2,2,0)&312&300&296&4\\
(2,1,1)&1260&1200&1188&12\\
(3,0,1),(3,1,0)&140&130&129&1\\
(4,0,0)&8&8&8&0.
\end{array}                                                   \tag{SL.5}
\]
Each of the two listed multidegrees has the displayed sizes.
The nullities sum to $180$.

The complete certificate is
\texttt{check\_genus\_three\_}\allowbreak\texttt{single\_line.py}.
For every block it computes a rank lower bound modulo~101 and
checks independent integer kernel lifts against the original
integer specialness matrix. The lower and upper bounds agree,
so these are rational ranks. Conjugation by $a$ is injective on
each degree-six multidegree used here: its free-Lie centralizer
consists of multiples of $a$. Indeed, the associative centralizer
of the free generator $a$ is $\Q[a]$: the coefficient equation
$aF=Fa$ along each finite chain of words with a non-$a$ core
forces those coefficients to vanish. The only primitive elements
of $\Q[a]$ in characteristic zero are the multiples of $a$,
as is seen from the binomial formula for its coproduct.
Hence $\Q[a]\cap\mathbb L_6=0$. Thus the selected column coordinates
describe the full derivation space without an unaccounted kernel.

On these 180 basis vectors the equations (SL.4) form an integer
matrix with $7977$ nonzero coordinate rows and rank $48$.
To certify the rank over $\Q$, the program selects 48 independent
rows modulo~101, performs exact rational row reduction on those
rows, clears denominators in all 132 free-coordinate kernel
vectors, and checks them against every original integer row.
Their independent free coordinates and the modular lower bound
prove that the full rational kernel has dimension $132$.
All these equations and ranks remain valid after extension to
every $\Q_\ell$.

\paragraph{The contraction image lies in two actual geometric brackets.}
Use right-nested brackets and define
\[
\begin{split}
 Z_1&=[\delta_{cdf},\delta_{cef},\delta_{aef},
                  \delta_{def},\delta_{bef},\delta_{cde}],\\
 Z_2&=[\delta_{bce},\delta_{adf},\delta_{adf},
                  \delta_{abd},\delta_{bce},\delta_{abc}].
\end{split}                                                   \tag{SL.6}
\]
Both belong to $J_6$ by the actual geometric Lie bracket.
Let $K$ be the vector of all 105 complete contractions of
the tensor associated with a special degree-six derivation.
The two rows indexed by
$(12)(34)(56)(78)$ and $(12)(36)(47)(58)$ give
\[
 \begin{pmatrix}-104&-288\\32&192\end{pmatrix},
             \qquad\det=-10752.                           \tag{SL.7}
\]
For every one of the 132 integer kernel vectors $D_j$, solve
for the two coefficients using (SL.7). The program checks
\emph{all} 105 coordinates of
\[
                 K(D_j)\in\operatorname{span}\{K(Z_1),K(Z_2)\}
                                                               \tag{SL.8}
\]
over the integers after multiplication by $-10752$.
The largest coefficient in the verified integer kernel is
$247830$, and all fixed-width multiplication bounds are checked.
The determinant is nonzero over every $\Q_\ell$; this step
makes no integral divisibility assertion.

Choose a basis of the one-dimensional $Q_6$ and regard $q_6$
as scalar valued. Symplectic complete reducibility extends it
equivariantly first to $\mathfrak h_{3,6}$ and then through
$\mathfrak h_{3,6}\hookrightarrow V^{\otimes8}$ to the full
tensor power. Every invariant functional on that tensor power
is a linear combination of the complete contractions, by
\cite[Proposition~2.1, p.~8]{AuditMoritaSakasaiSuzuki2015v2}.
Since it vanishes on $Z_1,Z_2$ by (SL.2), (SL.8) proves
$q_6(U_{3,6}\cap P_a)=0$.

\paragraph{The actual post-constraint projection.}
The diagonal torus preserves both $U_{3,6}$ and $P_a^{\rm all}$.
The quotient $Q_6$ has weight $t^3$. Take the algebraic weight
projector $e$ onto that weight in the finite-dimensional
derivation representation; it is a polynomial in a generic
diagonal torus element, with the denominators cleared in the
usual Lagrange formula. It satisfies
\[
 e(U_{3,6}\cap P_a^{\rm all})\subset U_{3,6}\cap P_a,
                    \qquad q_6(eD)=q_6(D).
\]
Applying the proved balanced assertion gives (SL.1).
The projection is applied to an already constrained vector in
the actual universal image. It is not asserted to be a new
individual Galois element.
\end{proof}

\paragraph{Arithmetic consequence on the original genus-three curve.}
Let $C/\Q$ be a smooth proper geometrically connected curve of
genus three, let $y\in C(\Q)$, and put $X=C\setminus\{y\}$.
For every prime $\ell$ and finite extension $K/\Q$, every nontrivial
closed procyclic subgroup of $\pi_1(X_{\overline K})^{(\ell)}$
whose arithmetic normalizer has open image in $G_K$ is contained
in an actual cusp inertia subgroup. No restriction on the
Jacobian is needed.

\begin{proof}
We apply (SL.1) to the original homology, so no CM decomposition
or permutation of character lines enters the proof.
Let $u$ generate the subgroup, $N$ be its arithmetic normalizer,
and $u_d$ its first nonzero rational Magnus term.
The geometric kernel of $N$ acts trivially on $u_d$, so the
normalizer exponent gives a character $\theta$ on the open
arithmetic image. The exact finite-cover detector of F.3.21
gives an actual covering $Y$ in whose rational homology $u$
has nonzero class, after finite arithmetic descent.
The intersection of the normalizer with that covering group
still has open arithmetic image; its geometric kernel acts
inner and trivially on the covering homology. Consequently
$u_d$ and this class have the same character $\theta$.

The original curve has one puncture, so its degree-one
homology is pure of weight one and its degree-$d$ free-Lie
part has weight $d$. If the detected class has nonzero image
in proper homology, a common good Frobenius gives $d=1$.
We exclude this case directly. Put $a=u_1\ne0$ and choose
a symplectic basis extending $a$. Over a finite extension
$L$ with $G_L\subset p(N)$, the actual finite-base-invariant
arithmetic quotient of Nakamura, pp.~115--117, supplies
\[
 \sigma\in G_L(6),\qquad D_\sigma\in U_{3,6},
                         \qquad q_6(D_\sigma)\ne0.
\]
Lift this same $\sigma$ to $N$. It acts identically on
homology, so its exponent on $u$ is one. If $X_u=\log u$
and $Z$ is the logarithm of the inner correction between
the normalizing lift and the canonical special representative,
the resulting equality is
\[
                   \phi_\sigma(X_u)=\exp(\operatorname{ad}Z)X_u.
\]
The left side first differs from $X_u$ in degree seven.
The degree-one term of $Z$ centralizes $a$ and is a rational
multiple of $a$. Multiplying the correction by a rational
power of the full $u$ removes that term without changing
the displayed conjugation. Successive degrees two through
five of $Z$ then centralize $a$; the associative-centralizer and
primitive-element argument above shows that the free Lie algebra
has no such homogeneous centralizer part. Comparing
degree seven gives
\[
                     D_\sigma(a)=[Z_6,a].
\]
This is exactly the nonprimitive-logarithm correction already
proved in F.3.22 and F.3.25, applied to one line. Now (SL.1)
contradicts the chosen nonzero quotient.

The detector's nonzero class therefore belongs to its boundary
kernel. After defining all boundary points, that kernel has
character $\chi_\ell$. Thus $\theta=\chi_\ell$ on the
intersected normalizer. Every element of the full normalizer
over that finite extension differs from such a lift by an
element of its geometric kernel, where $\theta$ is trivial.
The full normalizer is consequently cyclotomic with open
arithmetic image. Nakamura's original Theorem~2.1.1,
printed p.~94, applies to the class of finite $\ell$-groups
and gives actual cusp inertia containing the subgroup
\cite{AuditNakamura1994}. The argument concerns every prime.
\end{proof}

\paragraph{Separate all-genus balanced tangency lemma.}
For every genus $g\geq3$, a special degree-six derivation in
the $t^3$ weight that is tangent to a nonzero basis vector
$a$ is also tangent to its chosen symplectic partner $b$.
This assertion concerns the balanced derivation space; an
all-genus version of (SL.1) still needs its own proof.

\begin{proof}
Write $T=\sum_i(x_i\otimes D(y_i)-y_i\otimes D(x_i))$.
The specialness equation says $(1-\rho)T=0$, where $\rho$
moves the first tensor letter to the end. The tensor is thus
cyclic. In every balanced block its handle multiplicities
$n_i$ sum to four. If the constrained handle is absent,
both $D(a)$ and $D(b)$ vanish: their required multidegrees
are obtained by subtracting the partner letter from the
tensor multidegree and would otherwise have a negative entry.

If this handle occurs once, write $D(a)=[A,a]$; the
multidegree of $A$ contains neither $a$ nor $b$. The
$b$-leading part of $T$ is
\[
                      -bAa+baA.
\]
Each tensor word has a unique occurrence of $b$, so this
part determines its full cyclic orbit and hence determines
the whole block of $T$. Rotating to the unique $a$-leading
positions gives $-abA+aAb=a[A,b]$. Therefore
$D(b)=[A,b]$, as asserted, even for a block on four handles.

If the constrained handle has multiplicity at least two,
there are at most two other handles in its support. The
complete single-line block bases in (SL.5) are all
certified to be paired by the same final program.
For each polynomial $D(b)$ it groups words as
$b^i w b^j$, with $w$ beginning and ending in a letter
different from $b$, and solves the coefficient chains of
$[B,b]=D(b)$. It then verifies that exact commutator and
the Dynkin identity $\operatorname{Dyn}_6(B)=6B$, which
places $B$ in $\mathbb L_6$. These checks cover all 180
basis vectors, so they prove the actual inclusion of the
single-line space in the paired-line space over $\Q$ and
every $\Q_\ell$; the reverse inclusion is immediate.
Rename the at most three support
handles to obtain the required equality in every genus.
This argument gives the actual inclusion and equality of
the two specified balanced tangent spaces.
\end{proof}

\paragraph{Retained scope and source status.}
The full LM6 A.22 subgroup assertion still permits a rational
once-punctured detecting curve of arbitrary genus. Its genus-three
case is now handled without a simplicity or pairing hypothesis.
The all-genus mixed-normalizer calculation is not established
by the support bound alone. The MSS source is the opened
18 August 2026 v1 preprint. Its two separately recorded arithmetic
slips in the proofs on pp.~22 and 25 do not occur in (SL.3),
and the final argument above imports neither of the affected
scalar expansions.
